\chapter{CRONOGRAMA}
\label{cap:cronograma}

O cronograma foi atualizado para distinguir atividades concluídas, em andamento e planejadas. A revisão sistemática e sua documentação inicial foram concluídas. Em julho de 2026, o trabalho encontra-se na etapa de revisão metodológica e textual. O levantamento definitivo de tecnologias, bases de dados e alternativas de prototipação será realizado depois da consolidação dessas correções.

\begin{table}[htbp]
\centering
\caption{Cronograma atualizado do trabalho.}
\label{tab:cronograma}
\small
\begin{tabular}{p{3.0cm}p{8.0cm}p{2.5cm}}
\hline
\textbf{Período} & \textbf{Atividade} & \textbf{Status} \\
\hline
Março--abril/2025 & Definição do tema e revisão bibliográfica inicial & Concluído \\
Maio--novembro/2025 & Desenvolvimento do \textit{pipeline}, execução da revisão e elaboração do PTC & Concluído \\
Dezembro/2025--março/2026 & Consolidação dos dados, revisão dos estudos e ampliação do texto do TCC & Concluído \\
Abril--junho/2026 & Avaliação metodológica dos estudos e preparação dos artefatos da revisão & Concluído \\
Julho/2026 & Correções metodológicas, editoriais, normativas e de acessibilidade do texto & Em andamento \\
Após as correções & Levantamento e comparação de bases de dados, tecnologias e alternativas arquiteturais & Planejado \\
Etapa posterior & Implementação de modelos de referência e de um protótipo mínimo & Planejado \\
Etapa posterior & Avaliação técnica, análise das limitações e redação dos resultados & Planejado \\
Antes da submissão & Revisão final, normalização institucional e preparação para defesa & Planejado \\
\hline
\end{tabular}
\fonte{Elaborado pelo autor (2026).}
\end{table}

As datas das etapas posteriores dependerão da escolha da base, da disponibilidade dos dados e do desenho de avaliação aprovado. O cronograma não registra como concluídas atividades para as quais ainda não existem implementação, resultados ou documentação verificável.

%--------- FIM CRONOGRAMA ------------
